\documentclass[11pt]{article}
\usepackage[margin=1in]{geometry}
\usepackage{amsmath,amssymb,amsthm,mathtools,booktabs,array,microtype,hyperref}
\usepackage[T1]{fontenc}
\usepackage{lmodern}
\hypersetup{colorlinks=true,linkcolor=blue,citecolor=blue,urlcolor=blue}

\newtheorem{theorem}{Theorem}[section]
\newtheorem{lemma}[theorem]{Lemma}
\newtheorem{proposition}[theorem]{Proposition}
\newtheorem{corollary}[theorem]{Corollary}
\newtheorem{definition}[theorem]{Definition}
\newtheorem{remark}[theorem]{Remark}
\newcommand{\F}{\mathbb F}
\newcommand{\A}{\mathbb A}
\newcommand{\cB}{\mathcal B}
\newcommand{\cF}{\mathcal F}
\newcommand{\cR}{\mathcal R}
\newcommand{\OO}{\mathcal O}
\newcommand{\softO}{\widetilde O}
\newcommand{\eps}{\varepsilon}
\newcommand{\length}{\operatorname{length}}
\newcommand{\im}{\operatorname{im}}

\title{Prime-Field Summation-Polynomial Index Calculus:\\
An Entropy--Degree Audit and a Conditional-Class Barrier}
\author{Anonymous}
\date{26 November 2026}

\begin{document}
\maketitle

\begin{center}
\fbox{\parbox{0.92\linewidth}{\textbf{Research-draft status.}
This manuscript is not peer reviewed and its new claims are
\texttt{[pcc:unverified]}.  The barrier applies only to the explicitly defined
candidate/collision and quotient-materializing classes.  It does not prove a % [pcc:unverified]
lower bound for unrestricted arithmetic or rational circuits, or for ECDLP in
general.}}
\end{center}

\begin{abstract}
Let $E/\F_p$ be an ordinary, non-anomalous elliptic curve and let a prime-order
subgroup $G\subseteq E(\F_p)$ have order $n=\Theta(p)$.  We audit index-calculus
attacks that seek decompositions through Semaev polynomials with factor-base
$x$-set $F\subseteq\F_p$, $|F|=N$.  The audit separates three statements which
are often conflated.  First, the probability calculation is an exact averaging
upper bound; the usual asymptotic equality is a distribution heuristic.
Second, a finite faithful algebraic encoding of $F$ has scheme length at least
$N$, even when all its equations are quadratic.  Consequently $m$ independent
leaves have ambient quotient length at least $N^m$.  Third, this length is a
lower bound for algorithms that materialize the quotient, a multiplication
matrix, a Macaulay normal-form space, or all candidate sums, but it is not an
unconditional arithmetic-circuit lower bound for every possible solver.

For an explicitly defined class consisting of (i) generic candidate/collision
search, and (ii) faithful quotient-materializing Macaulay/resultant elimination,
we prove a square-root barrier.  In the stronger batched candidate model, for % [pcc:unverified]
fixed $m$ the optimal relation-collection exponent is at least
$m/(2m-1)>1/2$; as $m\to\infty$ it approaches $1/2$ from above.  We then audit
intervals, small-height rational reconstruction, multiplicative cosets,
cyclotomic and Dickson encodings, isogeny or torsion fibres, split $S_3$ trees,
and multipoint batching.  None escapes the stated class barrier.  We explicitly
identify the unproved escape hatch: a non-generic algorithm which exploits the
arithmetic circuit of a summation polynomial on a compact, positive-variable
encoding without enumerating or materializing its $N^m$ states.  Ruling that
out unconditionally would go substantially beyond current generic-group and
algebraic-elimination lower bounds.  Thus the result is a sharp barrier theorem
for a broad, auditable class, not a claimed unconditional lower bound for the
ECDLP itself.
\end{abstract}

\section{Setting and conventions}

Let
\[
 E: y^2=x^3+Ax+B,\qquad 4A^3+27B^2\ne0,
\]
over $\F_p$, where $p>3$.  We work in a cyclic subgroup $G=\langle P\rangle$
of prime order $n=\Theta(p)$ and are given $Q=qP$.  The qualifications
``non-anomalous'' and ``non-supersingular'' exclude familiar special attacks;
they do not by themselves imply a lower bound.

The $r$th Semaev polynomial $S_r$ is symmetric, absolutely irreducible in the
generic setting, and has degree $2^{r-2}$ in each coordinate.  In particular,
the equation for $m$ factor-base points and a fixed target is
\[
 S_{m+1}(X_1,\ldots,X_m,x_0)=0,
 \qquad \deg_{X_i}S_{m+1}=2^{m-1}.                 \tag{1}
\]
This corrects a common one-index shift: $2^{m-2}$ is the coordinate degree of
$S_m$, not $S_{m+1}$.  Nothing below depends on the convention; one may retain
the symbol $\delta_m=\deg_{X_i}S_{m+1}$.  A dense box treating the $m$
coordinate degrees independently has size
$\delta_m^m=2^{m(m-1)}$.  This is the source of the familiar
$2^{\Theta(m^2)}$ resultant/Macaulay estimate; it is not, by itself, a lower
bound for every possible solver.

Semaev introduced the prime-field bounded-coordinate approach
\cite{Semaev2004}; Gaudry and Diem developed the successful Weil-restriction
setting over extension fields \cite{Gaudry2009,Diem2011}.  Prime-field
algebraic encodings were studied by Petit--Kosters--Messeng
\cite{PetitKostersMesseng2016} and Amadori--Pintore--Sala
\cite{AmadoriPintoreSala2018}.  Existing lower-bound analyses for the naive
method use additional statistical hypotheses on Gr\"obner-basis behaviour
\cite{YokoyamaEtAl2020}.  Our claims below are deliberately different: they
are unconditional for the explicitly defined materialization/collision class,
and make no claim outside it.

Let $F\subseteq\F_p$ have size $N=p^{\alpha+o(1)}$.  For each $x\in F$ for
which $x^3+Ax+B$ is a square, choose either both points or a fixed sign.  Denote
the resulting point factor base by $\cB$ and put $t=|\cB|=\Theta(N)$.  Constants
of two and factors $m!$ never change a fixed-$m$ exponent, but we retain them
when they explain a probability.

\section{The exact cost and probability audit}

\subsection{Success probability: what is exact and what is heuristic}

Let $\Sigma_m$ be the set of admissible unordered multisets of $m$ factor-base
points (including whatever sign convention the algorithm uses), and let
\[
 \sigma_m:\Sigma_m\longrightarrow G,\qquad
 \{B_1,\ldots,B_m\}\longmapsto B_1+\cdots+B_m.
\]
For a uniformly random target $R=aP+bQ$, $R$ is exactly uniform in $G$ as soon
as either $a$ or $b$ is uniform.  Therefore
\[
 P_{\rm rel}=\Pr[R\in\im(\sigma_m)]
 =\frac{|\im(\sigma_m)|}{n}
 \le \min\left\{1,\frac{|\Sigma_m|}{n}\right\}
 \le \min\left\{1,\frac{\binom{t+m-1}{m}}{n}\right\}.               \tag{2}
\]
This is unconditional.  If repetitions are negligible and the sum map has
nearly Poisson fibres, then
\[
 P_{\rm rel}\sim 1-\exp\left(-\frac{t^m}{m!n}\right)
 \sim \frac{t^m}{m!n}\quad(t^m=o(n)),                               \tag{3}
\]
but both asymptotic equivalences are heuristics.  The barrier results below use
only (2), or exact expectation over independently uniform targets.

\subsection{The total cost equation}

The often written expression $P_{\rm rel}^{-1}T_{\rm solve}+T_{\rm LA}$ is the
cost of obtaining one successful target plus linear algebra; it omits the
number of independent rows required.  Let
\begin{itemize}
 \item $M\ge t+1$ be the number of accepted rows (including oversampling),
 \item $\rho\le1$ be the conditional probability that an accepted row raises
       the rank,
 \item $h\ge1$ be the expected number of usable rows returned by one successful
       solve, and
 \item $T_{\rm trial}$ include target generation, system construction, solving,
       rationality filtering, and verification.
\end{itemize}
Then the auditable classical cost is
\[
 \boxed{\quad
 T_{\rm total}=T_{\rm FB}+T_{\rm pre}
 +\frac{M}{\rho hP_{\rm rel}}T_{\rm trial}
 +T_{\rm LA}(t,M).\quad}                                             \tag{4}
\]
For a solver returning at most one row, $h=1$.  With a sparse $M\times t$
matrix having $O(m)$ nonzeros per row, Wiedemann/Lanczos gives the standard
upper estimate
\[
 T_{\rm LA}=\softO(mt^2)=p^{2\alpha+o(1)}                            \tag{5}
\]
field operations; dense elimination gives $O(t^\omega)$.  Equation (5) is an
algorithmic cost model, not a proven lower bound for arbitrary structured
linear algebra.  At least $t$ independent scalar unknowns must nevertheless be
resolved in the traditional fixed-factor-base formulation.

Using the heuristic sparse-fibre value (3), (4) becomes
\[
 T_{\rm rel}=p^{1-(m-1)\alpha+o(1)}T_{\rm trial}                    \tag{6}
\]
for $M=\Theta(N)$ and $h,\rho=p^{o(1)}$.  Thus a claimed exponent must include
the exponent of $T_{\rm trial}$; setting it to $p^{o(1)}$ is exactly the
unjustified step in most prime-field proposals.

\section{Faithful encodings and the invariant that degree alone misses}

\begin{definition}[Faithful finite encoding]
A faithful finite encoding of $F$ is a zero-dimensional affine $\F_p$-scheme
$Z\subseteq\A^k$, together with a rational map $\pi:Z\dashrightarrow\A^1$
defined at all relevant rational points, such that $\pi(Z(\F_p))=F$.  Extension
points and extra rational witnesses are permitted, so this definition is
favourable to an attacker.  A solver is \emph{quotient-materializing} if it
constructs a basis of $\F_p[Z]$, the product algebra $\F_p[Z]^{\otimes m}$, a
multiplication matrix on it, or a Macaulay/resultant object containing an
equivalent normal-form space.
\end{definition}

\begin{lemma}[Entropy--length lemma]\label{lem:length}
Every faithful finite encoding satisfies
\[
 \deg Z=\length_{\overline{\F}_p}(Z)\ge |Z(\F_p)|\ge |F|=N.          \tag{7}
\]
For $m$ independently selected leaves,
\[
 \length(Z^m)=(\length Z)^m\ge N^m.                                 \tag{8}
\]
Consequently every quotient-materializing solver performs
$\Omega(N^m)$ scalar writes or operations merely to materialize such a basis
or matrix domain.
\end{lemma}

\begin{proof}
The length of a zero-dimensional scheme over an algebraic closure is the sum
of its local multiplicities and is at least its number of distinct rational
points.  Surjectivity gives at least one rational witness for every element of
$F$.  Length is multiplicative under products.  The last sentence follows
from the definition of quotient-materializing.
\end{proof}

This lemma deliberately does \emph{not} say that one equation has degree $N$.
For example, if $N=2^k<p$, the interval $\{0,\ldots,N-1\}$ has the quadratic
encoding
\[
 b_j^2-b_j=0\ (0\le j<k),\qquad
 x=\sum_{j=0}^{k-1}2^jb_j.                                          \tag{9}
\]
All equations have degree two, but the Boolean scheme has length $2^k=N$.
For $m$ leaves its product quotient has the $N^m$ square-free monomials
$\prod b_{ij}^{e_{ij}}$ as a basis.  A total-degree Macaulay construction can
be still larger; a perfectly multigraded implementation reaches, but does not
beat, this entropy dimension.  Circuit lifting can similarly replace any % [pcc:unverified]
high-degree computation by quadratics while increasing the number of
variables.  Hence maximum equation degree is not an invariant; scheme length
and the relevant Hilbert function are.

For a univariate exact encoding the familiar statement remains true.

\begin{lemma}[Univariate and fibre degree]
If a nonzero $g\in\F_p[X]$ vanishes on $F$, then $\deg g\ge N$.  If $F$ is
contained in one fibre of a rational map $\phi:\mathbb P^1\dashrightarrow
\mathbb P^1$ of degree $d$, away from poles, then $N\le d$.  More generally,
the union of $r$ fibres has at most $rd$ geometric points counted without
multiplicity.
\end{lemma}

\begin{proof}
The first assertion is the root bound.  A fibre of a degree-$d$ map has scheme
length $d$; taking a union gives the last assertion.
\end{proof}

Thus $X^p-X$ contributes degree $p$ and does not help.  Replacing it by
$g_F(X)=\prod_{a\in F}(X-a)$ contributes degree $N$.  Multiple bounded-degree
variables can reduce individual degrees, as (9) shows, but cannot reduce the
finite scheme length.

\section{A precise barrier class}

No unconditional $p^{1/2-o(1)}$ lower bound is known for arbitrary algorithms
using the coordinates of a general elliptic curve.  Such a result would be a
major ECDLP lower bound.  We therefore state exactly the class covered.

\begin{definition}[Audited prime-field summation class $\mathsf{APS}$]
An attack is in $\mathsf{APS}$ if it uses a fixed factor base of size
$t=\Theta(N)$ and every relation-finding stage is a composition of the
following operations.
\begin{enumerate}
 \item A faithful finite encoding of each factor-base leaf, with direct
       $S_{m+1}$, a split addition tree, or any bounded-fibre algebraic
       correspondence whose exceptional components are saturated away.
 \item Either (E) explicit generation/materialization of candidate group sums
       and target group elements followed by equality, hashing, sorting, or
       multipoint comparison; or (Q) quotient-materializing Macaulay,
       resultant, multiplication-matrix, or normal-form elimination as in
       Lemma~\ref{lem:length}.
 \item Precomputation is charged once, target batching is allowed, and every
       generated relation and every output coefficient is charged.  Generic
       group operations may be arranged arbitrarily.
\end{enumerate}
Algorithms that exploit a special coordinate circuit without materializing
candidate states or the finite quotient are, intentionally, outside the
class.  This boundary is necessary rather than cosmetic.
\end{definition}

\subsection{The batched collision theorem}

Suppose a batch materializes $D$ distinct unordered candidate tuples, each
carrying its group sum, and $B$ distinct, uniform target elements (their
coefficients in $P,Q$ are retained).  Candidate sums may coincide: candidate
tuples are counted with multiplicity precisely so that several decompositions
of one group element are not lost.  Conditional on the tuples, the exact
expected number of candidate-tuple--target matches is at most $DB/n$.
Target--target collisions can directly solve the DLP, but in the generic group
model they already require square-root work and do not create an exception.

\begin{theorem}[Batched candidate barrier]\label{thm:batch}
In mode (E), collecting $N$ candidate--target matches has expected materialized
work
\[
 C\ge 2\sqrt{Nn}=p^{(1+\alpha)/2+o(1)}.                              \tag{10}
\]
If at most
\[
 K\le\binom{t+m-1}{m}=p^{m\alpha+o(1)}                              \tag{11}
\]
different unordered $m$-leaf candidates exist, the sharper exponent is
\[
 c_m(\alpha)=
 \begin{cases}
 1-(m-1)\alpha,& \alpha<1/(2m-1),\\[2mm]
 (1+\alpha)/2,& \alpha\ge1/(2m-1).
 \end{cases}                                                        \tag{12}
\]
For fixed $m$, optimizing the factor-base exponent gives
\[
 \min_{\alpha\ge0}c_m(\alpha)=\frac{m}{2m-1}>\frac12,\qquad
 \alpha_{\rm opt}=\frac1{2m-1}.                                   \tag{13}
\]
The bound tends to $1/2$ from above as $m\to\infty$.
\end{theorem}

\begin{proof}
For batches indexed by $j$, exact averaging gives at most
$\sum_jD_jB_j/n$ expected matches.  If
$C=\sum_j(D_j+B_j)$, then
$\sum_jD_jB_j\le C^2/4$ (concentrating the budget in one balanced batch only
increases the product).  Requiring expected output $N$ yields (10).

In one optimal batch one must have $DB\ge Nn$ and $D\le K$.  If
$K\ge\sqrt{Nn}$, the balanced choice has cost exponent $(1+\alpha)/2$.  If
$K<\sqrt{Nn}$, the best possible choice $D=K$, $B=Nn/K$ is dominated by
$B=p^{1-(m-1)\alpha+o(1)}$.  The boundary is
$m\alpha=(1+\alpha)/2$, namely $\alpha=1/(2m-1)$; both branches there equal
$m/(2m-1)$.
\end{proof}

This theorem already grants arbitrary batching and shared evaluation.  It is
therefore stronger than analysing one target at a time.  It is also compatible
with the Nechaev--Shoup generic-group lower bound: an entirely generic ECDLP
algorithm needs $\Omega(\sqrt n)$ group operations.  The extra factor
$N^{1/2}$ in (10) comes from demanding $N$ traditional factor-base rows.

\begin{theorem}[Barrier for $\mathsf{APS}$]\label{thm:main}
Let $n=\Theta(p)$ and $N=p^{\alpha+o(1)}$.
\begin{enumerate}
 \item Every mode-(E) attack has relation-collection exponent at least
       $c_m(\alpha)$ in (12), and at least $m/(2m-1)>1/2$ after optimizing
       $\alpha$ for fixed $m$.
 \item Every mode-(Q) attack materializing the independent-leaf quotient has
       an ambient algebraic state space of dimension at least $N^m$.  A fresh
       per-target elimination is already that expensive.  Even granting a
       reusable table/eliminant and constant-time query and recovery, its cost
       plus the number of random target queries needed for $N$ rows is at least
       \[
          N^m+\frac{Nn}{N^m}
          =p^{m\alpha+o(1)}+p^{1-(m-1)\alpha+o(1)},
       \]
       whose optimized exponent is $m/(2m-1)$.  Splitting, quadratic circuit
       lifting, sparsity of the leaf equations, and variable changes do not
       alter this length.
 \item Allowing $m\to\infty$ can make the mode-(E) exponent
       $1/2+o(1)$ but never $1/2-\eps$.  Factors $m!$, torsion symmetries, and
       bounded-degree correspondences only change the $p^{o(1)}$ term when
       $m=p^{o(1)}$ and the multiset count is used exactly.
\end{enumerate}
Hence no attack in $\mathsf{APS}$ has total cost $p^{1/2-\eps}$ for fixed
$\eps>0$.  In particular it has no $L_p[1,c]$ complexity with $c<1/2$.
\end{theorem}

\begin{proof}
Part 1 is Theorem~\ref{thm:batch}; adding factor-base construction and linear
algebra cannot lower a cost.  For Part 2, Lemma~\ref{lem:length} charges the
materialized quotient.  Equation (2) bounds success by $N^m/n$ up to
$p^{o(1)}$ symmetry factors, so even free post-precomputation queries require
at least $Nn/N^m$ trials for $N$ rows.  Balancing the two terms gives
$m\alpha=1-(m-1)\alpha$, hence $\alpha=1/(2m-1)$ and exponent
$m/(2m-1)$.  Part 3 follows from the same exponent.  The final claim follows by
taking the minimum over the two permitted modes, with a direct generic
target--target collision covered by the $\Omega(\sqrt n)$ bound
\cite{Shoup1997}.
\end{proof}

The theorem is unconditional \emph{inside its stated computational class}.
It is not a lower bound for a solver outside that class, and it does not use a
semi-regularity or first-fall-degree assumption.

\section{Audit of the proposed structural instruments}

\subsection{Direct substitution and lifting}

Let $x=\phi(u_1,\ldots,u_k)$ be a polynomial map of degree $d$.  Direct
substitution changes the coordinate degree $\delta_m$ of (1) to at most
$d\delta_m$ in the corresponding parameter block (and generically attains
that product).  A rational map $x=a(u)/b(u)$ has the same effect after clearing
denominators.  Introducing $x$ and the equation $b(u)x-a(u)=0$ avoids this
syntactic multiplication, but it does not change the finite-domain length.
If the parameters range over finite sets $U_j$, exact algebraic constraints
have product-scheme length $\prod_j|U_j|$, at least $N$ divided by the maximal
fibre size of $\phi$.  Thus substitution pays in degree; faithful lifting pays
in variables and quotient length.

If the $u_j$ range freely over $\F_p$, the equations $u_j^p-u_j$ have degree
$p$, and a nonconstant bounded-degree map normally has an image of order $p$,
not $p^\alpha$.  If the parameter variety is left positive-dimensional over
$\overline{\F}_p$, it is not a finite factor base.  Solving over the algebraic
closure then admits extension-field points which must be rejected; imposing
$\F_p$-rationality restores the degree-$p$ field equations or an equivalent
finite encoding.

\subsection{Integer intervals and Boolean encodings}

For $F=\{0,\ldots,N-1\}$, the univariate vanishing polynomial has degree $N$.
Substitution by an integer parameter does not make ``smallness'' algebraic:
without a range constraint one solves over all of $\F_p$.  The Boolean encoding % [pcc:unverified]
(9) is the best response to this objection: it has $k=\lceil\log_2N\rceil$
quadratic constraints and a linear map to $x$.  It genuinely decouples maximum
equation degree from density.

However, for $m$ leaves the Boolean quotient has dimension $2^{km}\ge N^m$.
Branching enumerates those states, a meet-in-the-middle is a candidate method,
and full Macaulay/normal-form elimination materializes the quotient.  All fall
under Theorem~\ref{thm:main}.  A new solver exploiting the particular circuit
$S_{m+1}(\sum2^jb_{1j},\ldots,x_0)$ without doing any of these things is outside
the theorem and is the clearest remaining research target.  No regularity
claim about that Boolean system is made here.

\subsection{Small-height rational reconstruction}

Take
\[
 x=u/v\pmod p,\qquad 1\le |u|,|v|\le L,\quad v\ne0.               \tag{14}
\]
For $L^2<p/2$, congruent reduced fractions are equal over the integers, and the
number of distinct ratios is $\Theta(L^2)$ (up to the usual coprimality
constant).  Thus $N=\Theta(L^2)$.  Exact interval constraints for $u$ and $v$
have length $\Theta(L^2)$ per leaf.  After clearing denominators, the Semaev
equation need not have degree $L$, but the product quotient has length
$L^{2m}=\Theta(N^m)$.  Bit-decomposing $u,v$ merely converts this to the Boolean
case.  If $L\ge p^{1/2}$, modular collisions become abundant and the image is
already a positive fraction of $\F_p$, so it is not a small factor base.

\subsection{Multiplicative subgroups, cosets, cyclotomic and Dickson sets}

If $H\le\F_p^*$ has order $N$, a coset is cut out by the sparse polynomial
$X^N-c$.  Sparsity makes membership testing fast, but its zero-dimensional
coordinate algebra still has dimension $N$; $m$ copies have dimension $N^m$.
An FFT over $H^m$ evaluates all states in essentially $N^m$ time, while partial
sum tables are mode (E).

There is no hidden algebraic homomorphism converting multiplicative convolution
to addition on a general elliptic curve.  Indeed a rational map
$\mathbb G_m\dashrightarrow E$ extends, because $E$ is proper, to a morphism
$\mathbb P^1\to E$, and Riemann--Hurwitz forces it to be constant.  The same is
true for $\mathbb G_a$.  Over the prime field, moreover, the only additive
subgroups of $(\F_p,+)$ are $\{0\}$ and $\F_p$.  Cyclotomic and Dickson
polynomials can compress the description or expose a $\mathbb G_m$ quotient;
they do not supply the missing homomorphism.  Their root count is still bounded
by degree and their faithful finite scheme still has length at least $N$.

This does not prove that every sparse-polynomial algorithm is slow.  It proves % [pcc:unverified]
that sparse membership alone supplies no entropy reduction; a successful
algorithm would have to exploit a further, non-generic interaction with the
elliptic addition circuit.

\subsection{Torsion shifts, isogenies, and modular correspondences}

Translation by a fixed torsion point and symmetrization can reduce constants,
monomial counts, or a finite symmetry factor.  A fixed-degree isogeny
$\psi:E'\to E$ is a group homomorphism, but a fibre has length $\deg\psi$.
Obtaining $N$ factor-base elements from $r$ selected fibres requires
$r\deg\psi\ge N$.  If the isogeny degree grows like $N$, its rational maps and
division polynomials carry the same growing degree; if it is bounded, the
entropy remains in the number of selected fibres.  A bounded-bidegree
correspondence has the identical fibre-degree obstruction.  Exceptional curves
with special endomorphisms are outside the stated ``general curve'' target and
must be analysed separately; bounded automorphism groups only change constants.

\subsection{Split trees of S3 equations}

For reference,
\begin{align*}
S_3(X_1,X_2,Z)={}&(X_1-X_2)^2Z^2
-2\big((X_1+X_2)(X_1X_2+A)+2B\big)Z\\
&+(X_1X_2-A)^2-4B(X_1+X_2),
\end{align*}
which is quadratic in each variable.  A chain or balanced tree replaces (1)
by $m-1$ copies of $S_3$ and $m-2$ internal $x$-variables (with the fixed target
at the root).  Before factor constraints, it has
\[
 2m-2\text{ variables},\quad m-1\text{ equations},\quad
 \dim=m-1,                                                         \tag{15}
\]
the same dimension as one equation in the $m$ leaf variables.  The lift has
distributed the eliminant degree; it has not made the leaf choices finite.

After attaching faithful leaf encodings, the ambient leaf algebra still has
length at least $N^m$.  For each nonsingular leaf tuple the internal variables
have only $2^{O(m)}$ sign witnesses.  Saturating addition denominators and the
vertical/2-torsion exceptional loci removes parasitic components but cannot
reduce leaf entropy.  Conversely, failing to saturate can only add spurious
components and optimistic timings.

There is a rationality subtlety.  $S_r=0$ characterizes points over
$\overline{\F}_p$.  Leaf $x$-coordinates must first be certified to lift to
$E(\F_p)$, and internal roots returned in extensions must be rejected unless
one uses explicit rational addition variables.  Adding $Z^p-Z$ certifies an
internal $x$ but reintroduces degree $p$.  This filtering never raises the
success probability and therefore cannot improve an exponent. % [pcc:unverified]

\subsection{Target batching and multipoint resultants}

Suppose $D$ leaf tuples are evaluated once and compared with $B$ targets.
Multipoint evaluation, hashing, and product trees can reduce logarithmic
factors, but exact averaging still gives at most $DB/n$ expected matches.
Theorem~\ref{thm:batch} already optimizes $D+B$ subject to producing $N$
matches.  An alleged sublinear-in-$B$ batch solver must at least avoid
materializing its $B$ target coefficients or its outputs; if it produces
$\Theta(B)$ distinct rows, output size alone is $\Omega(B)$.

On the elimination side, the norm/resultant of the relation function over
$Z^m$ is a determinant on an algebra of dimension at least $N^m$.  Fast
multipoint evaluation amortizes evaluating a precomputed eliminant but does not
make that eliminant free.  A method that represents and interrogates this norm
by a smaller special-purpose circuit without materialization is, again, the
explicitly open escape hatch rather than an audited speedup. % [pcc:unverified]

\section{Macaulay, Hilbert, and zero-dimensionality audit}

For a box encoding with variables $u_1,\ldots,u_k$ and constraints
$g_j(u_j)=0$, $\deg g_j=L_j$, the pre-relation quotient has the normal monomial
basis
\[
 u_1^{e_1}\cdots u_k^{e_k},\qquad 0\le e_j<L_j,                    \tag{16}
\]
of size $\prod_jL_j$.  Across $m$ leaves the multigraded Hilbert function
therefore stabilizes at $(\prod_jL_j)^m\ge N^m$.  A dense total-degree Macaulay
matrix at degree $D$ has
\[
 \binom{r+D}{r}                                                     \tag{17}
\]
columns for $r$ variables, often much larger than the multigraded bound.  We do
not assert a semi-regular formula for the degree of regularity: split Semaev
systems are special, and first-fall degree is not a proved proxy for regularity. % [pcc:unverified]

The same invariant appears in classical elimination.  Let $A$ be the product
coordinate algebra of the $m$ leaf encodings and let
$d_A=\dim_{\F_p}A\ge N^m$.  For a target-dependent relation polynomial $f_T$,
the norm eliminant is
\[
 \operatorname{Norm}_{A/\F_p}(f_T)=\det(m_{f_T}),                  \tag{18}
\]
where $m_{f_T}$ is a $d_A\times d_A$ multiplication matrix.  A Macaulay or
resultant implementation that constructs this normal-form action therefore
has at least $d_A$ independent basis columns (usually many more before
reduction).  A special-purpose circuit representing and querying the same norm
without materializing this space is outside mode (Q), exactly as required by
the stated boundary of the theorem.

A proposed system must pass all of the following tests.
\begin{enumerate}
 \item \textbf{Dimension.}  Without finite leaf constraints, the direct or
 split relation variety has dimension $m-1$.  A positive-dimensional solver
 output is not a factor-base decomposition algorithm.
 \item \textbf{Rationality.}  Solutions over $\overline{\F}_p$ are not
 automatically in $\F_p$.  Field equations, certified parameter domains, or
 post-verification must be charged.
 \item \textbf{Saturation.}  Denominators, vertical additions, repeated leaves,
 and torsion exceptional loci must be saturated or handled separately.
 \item \textbf{Symmetry.}  Permuting $m$ leaves gives at most an $m!$ factor.
 Quotienting this symmetry changes fixed-$m$ constants, not the exponent.
 \item \textbf{Rank.}  Counting polynomial zeros is not counting independent
 relation rows.  The factors $\rho$ and $h$ in (4) must be measured or bounded.
\end{enumerate}

\section{Small-parameter sanity checks}

Take first the common constant-success choice $N=p^{1/m+o(1)}$.  Then
$N^m\asymp p$, so (3) predicts constant success, but a quotient-materializing
solver already sees $p^{1+o(1)}$ ambient states.  The strongest explicit
candidate batching lower bound is $p^{(m+1)/(2m)+o(1)}$.  If $N$ is optimized
instead, Theorem~\ref{thm:batch} gives $N=p^{1/(2m-1)}$ and exponent
$m/(2m-1)$.

\begin{center}
\small
\begin{tabular}{c c c c c c}
\toprule
$m$ & polynomial & coordinate degree & constant-success $\alpha$
& batched exponent & optimized exponent\\
\midrule
3 & $S_4$ & $4$ & $1/3$ & $2/3$ & $3/5$\\
4 & $S_5$ & $8$ & $1/4$ & $5/8$ & $4/7$\\
5 & $S_6$ & $16$& $1/5$ & $3/5$ & $5/9$\\
\bottomrule
\end{tabular}
\end{center}

For $m=3$, a split target equation uses two $S_3$ nodes and one internal
variable.  With $N=p^{1/3}$, its leaf quotient has length $N^3=p$; explicit
batching needs at least $p^{2/3}$ work.  Optimizing at $N=p^{1/5}$ balances
$D=N^3$ candidates and $B=p/N^2$ targets, both $p^{3/5}$.

For $m=4$, a balanced target tree may use
\[
 S_3(x_1,x_2,z_1)=S_3(x_3,x_4,z_2)=S_3(z_1,z_2,x_0)=0.
\]
It has six variables and three equations before leaf constraints, hence
dimension three.  At $N=p^{1/4}$ the quotient length is $p$ and the candidate
bound is $p^{5/8}$.  The optimized choice $N=p^{1/7}$ gives $p^{4/7}$.

For $m=5$, any chain or balanced tree has eight variables and four $S_3$
equations before constraints, hence dimension four.  At $N=p^{1/5}$ the
candidate bound is $p^{3/5}$; optimizing at $N=p^{1/9}$ gives $p^{5/9}$.
In every row the sparse-linear-algebra estimate $p^{2\alpha}$ is below the
optimized collection exponent, so it does not rescue or invalidate the
barrier.

\section{What has and has not been proved} % [pcc:unverified]

The following conclusions require no regularity heuristic.
\begin{enumerate}
 \item The exact average success upper bound is (2), and the complete cost is
       (4), not a one-relation expression.
 \item Every faithful finite algebraic encoding has length at least $N$ per
       leaf.  Low individual degree is possible only by paying in variables;
       the $m$-leaf product length remains at least $N^m$.
 \item Split trees preserve the positive dimension before leaf constraints and
       preserve leaf entropy afterwards.
 \item Explicit candidate generation, even with arbitrary target batching,
       has fixed-$m$ optimum exponent $m/(2m-1)>1/2$.
 \item Intervals, rational reconstruction, subgroup/cyclotomic/Dickson sets,
       bounded-degree fibres, torsion symmetries, and multipoint evaluation do
       not escape these invariants.
\end{enumerate}

What is \emph{not} proved is that every arithmetic circuit solving a faithfully % [pcc:unverified]
encoded summation system must enumerate candidates or materialize the quotient.
Proving such a statement would require new lower bounds for a highly structured
family of finite-field polynomial systems and, if made broad enough, would
subsume a non-generic ECDLP lower bound.  Likewise, no claim is made that a
first-fall degree equals a degree of regularity, or that Semaev restrictions
behave as random Boolean functions.

Accordingly, the answer to the successor target is (B), with an explicit class:
there is a sharp $p^{1/2}$ barrier for generic candidate/collision attacks and
faithful quotient-materializing algebraic elimination, and every proposed
mechanism in the question falls into that class unless it supplies a genuinely
new non-materializing circuit algorithm.  The most concrete research problem
left open is therefore:
\begin{quote}
Given the Boolean or small-height circuit encoding of $F$, can one solve the
one-parameter family of restricted $S_{m+1}$ equations in
$p^{1/2-\eps}$ total time without enumerating candidate sums, materializing a
$N^m$ quotient, or using an unproved regularity assumption?
\end{quote}
No presently justified mechanism in the audited families does so.

\begin{thebibliography}{99}
\bibitem{Semaev2004}
I. Semaev,
\emph{Summation polynomials and the discrete logarithm problem on elliptic curves},
IACR ePrint 2004/031, 2004.
\url{https://eprint.iacr.org/2004/031}

\bibitem{Diem2011}
C. Diem,
\emph{On the discrete logarithm problem in elliptic curves},
Compositio Mathematica 147 (2011), 75--104.
\url{https://doi.org/10.1112/S0010437X10005075}

\bibitem{Gaudry2009}
P. Gaudry,
\emph{Index calculus for abelian varieties of small dimension and the elliptic
curve discrete logarithm problem},
Journal of Symbolic Computation 44 (2009), 1690--1702.
\url{https://doi.org/10.1016/j.jsc.2008.08.005}

\bibitem{PetitKostersMesseng2016}
C. Petit, M. Kosters, and A. Messeng,
\emph{Algebraic approaches for the elliptic curve discrete logarithm problem
over prime fields}, PKC 2016, LNCS 9615, 3--18.
\url{https://doi.org/10.1007/978-3-662-49387-8_1}

\bibitem{AmadoriPintoreSala2018}
A. Amadori, F. Pintore, and M. Sala,
\emph{On the discrete logarithm problem for prime-field elliptic curves},
Finite Fields and Their Applications 51 (2018), 168--182.
\url{https://doi.org/10.1016/j.ffa.2018.01.009}

\bibitem{YokoyamaEtAl2020}
K. Yokoyama, M. Yasuda, Y. Takahashi, and J. Kogure,
\emph{Complexity bounds on Semaev's naive index calculus method for ECDLP},
Journal of Mathematical Cryptology 14 (2020), 460--485.
\url{https://doi.org/10.1515/jmc-2019-0029}

\bibitem{HuangKostersYeo2015}
M.-D. A. Huang, M. Kosters, and S. L. Yeo,
\emph{Notes on summation polynomials}, arXiv:1503.08001.
\url{https://arxiv.org/abs/1503.08001}

\bibitem{Shoup1997}
V. Shoup,
\emph{Lower bounds for discrete logarithms and related problems},
EUROCRYPT 1997, LNCS 1233, 256--266.
\url{https://www.shoup.net/papers/dlbounds1.pdf}

\bibitem{FaugereEtAl2014}
J.-C. Faug\`ere, L. Huot, A. Joux, G. Renault, and V. Vitse,
\emph{Symmetrized summation polynomials: using small order torsion points to
speed up elliptic curve index calculus}, EUROCRYPT 2014, LNCS 8441, 40--57.
\url{https://www.iacr.org/archive/eurocrypt2014/84410158/84410158.pdf}
\end{thebibliography}

\end{document}
